% =============================================================================
%  postextual/anexos.tex
%  ANEXOS — documentos de TERCEIROS incluídos para embasar o trabalho
%
%  Exemplos: trechos de documentação oficial de APIs, normas técnicas,
%            especificações de protocolos, contratos, documentos
%            institucionais fornecidos por terceiros.
%
%  Diferença importante:
%    Apêndice → criado por VOCÊ
%    Anexo    → criado por TERCEIROS
%
%  Numeração: Anexo A, Anexo B, Anexo C...
% =============================================================================

\begin{anexosenv}

% \partanexos gera a página divisória com o título "ANEXOS" em destaque.
\partanexos


% -----------------------------------------------------------------------------
%  ANEXO A
%  Substitua pelo seu anexo real.
% -----------------------------------------------------------------------------
\chapter{TODO: Título do Anexo A}
\label{anx:doc-externo}

TODO: Descreva brevemente a origem deste documento e como ele se
relaciona com o trabalho desenvolvido.

% Opção 1: incluir conteúdo textual diretamente
TODO: Cole ou transcreva aqui o trecho relevante do documento externo,
indicando a fonte de origem.

% Opção 2: incluir um PDF externo inteiro (requer \usepackage{pdfpages} no main.tex)
% \includepdf[pages=-]{imagens/nome_do_arquivo_externo.pdf}
% O parâmetro pages=- inclui TODAS as páginas do PDF.
% Para incluir apenas páginas específicas: pages={1,3,5} ou pages={2-4}


% -----------------------------------------------------------------------------
%  ANEXO B (descomente e adapte se necessário)
% -----------------------------------------------------------------------------
% \chapter{TODO: Título do Anexo B}
% \label{anx:norma}
%
% TODO: Conteúdo do Anexo B.


\end{anexosenv}
